\documentclass[11pt]{article}
\usepackage[a4paper,margin=1in]{geometry}
\usepackage{amsmath,amssymb,amsthm}
\usepackage{graphicx}
\usepackage[colorlinks=true,linkcolor=blue,citecolor=blue,urlcolor=blue]{hyperref}

\newtheorem{theorem}{Theorem}
\newtheorem{lemma}[theorem]{Lemma}
\newtheorem{remark}[theorem]{Remark}
\newcommand{\N}{\mathbb N}
\newcommand{\E}{\mathcal E}

\title{No Uniform Polynomial Exponent Gap for Positive\\
Kreiss Bounded Semigroups on $L^2$}
\author{Run \texttt{fa\_banach\_001}, lane 0}
\date{22 July 2026}

\begin{document}
\maketitle

\begin{center}
\textbf{Status: partial result, likely valid; human review requested.}
\end{center}

\section{Source question}

Arnold and Coine ask in Question 3.9(2) of \cite{ArnoldCoine} whether every
uniformly eventually positive Kreiss bounded $C_0$-semigroup $(T_t)_{t\geq0}$
on an $L^p$-space, $1<p<\infty$, satisfies
\[
  \|T_t\|=O(t^{1-\varepsilon})
\]
for some $\varepsilon\in(0,1)$.  The question appears on page 11 of the
arXiv PDF and is reproduced below.

\begin{center}
\includegraphics[width=\textwidth]{figures/open_problem_crop.png}
\end{center}

We do not settle the quantifier ``for each semigroup there exists
$\varepsilon$.''  We prove a sharpness result: $\varepsilon$ cannot be chosen
as a function of the ambient exponent $p$ alone, even at $p=2$ and even for
semigroups that are positive from time zero.

\section{Proof intuition}

Bonilla and Muller constructed, for every prescribed $\delta>0$, a positive
Cesaro-bounded operator $T$ on a Hilbert lattice whose powers are bounded below
by a constant times $n^{1-\delta}$ \cite[Theorem 2.3]{BonillaMuller}.  We turn
the discrete dynamics into a continuous semigroup by moving functions around
the unit interval and applying one copy of $T$ each time the moving point
crosses the endpoint.

At integer times this suspension acts pointwise as $T^n$, so it preserves the
large power norms exactly.  The only nonformal step is to show that Cesaro
boundedness also survives suspension.  Over an integer interval $[0,N]$, the
integral of the suspension splits into an initial fragment, the discrete
partial sum $I+T+\cdots+T^{N-1}$ applied to the integral of the function, and
a final fragment multiplied by $T^N$.  The first two are controlled directly,
and $T^N$ is the difference of two consecutive partial sums.  Hence the
continuous Cesaro means are uniformly bounded.  Positivity then lets us invoke
the Cesaro--Kreiss equivalence from \cite[Proposition 2.6]{ArnoldCoine}.

\section{The suspension lemma}

Let $X$ be a Banach space and $T\in\mathcal L(X)$.  On
\[
  \E:=L^2([0,1);X)
\]
define, for $t\geq0$ and $0\leq s<1$,
\begin{equation}\label{eq:suspension}
  (S(t)f)(s):=T^{m(t,s)}f\bigl(s+t-m(t,s)\bigr),
  \qquad m(t,s):=\lfloor s+t\rfloor.
\end{equation}

\begin{lemma}[Basic suspension properties]\label{lem:basic}
The family $(S(t))_{t\geq0}$ is a $C_0$-semigroup on $\E$.  If $X$ is a
Banach lattice and $T$ is positive, then every $S(t)$ is positive.  Moreover,
for every $n\in\N_0$,
\begin{equation}\label{eq:integer}
  \|S(n)\|=\|T^n\|.
\end{equation}
\end{lemma}

\begin{proof}
Writing $s+t=m+r$ with $m\in\N_0$ and $r\in[0,1)$, and then writing
$r+u=m'+r'$, gives
\[
  m(t+u,s)=m+m',\qquad
  s+t+u-m(t+u,s)=r'.
\]
This proves $S(t+u)=S(t)S(u)$.  For fixed $t$, the integer $m(t,s)$ takes at
most two values, so $S(t)$ is bounded.

For $0<t<1$ one has
\[
 (S(t)f)(s)=
 \begin{cases}
   f(s+t),&0\leq s<1-t,\\
   Tf(s+t-1),&1-t\leq s<1.
 \end{cases}
\]
Translation continuity in $L^2$, together with absolute continuity of the
integral on the two boundary intervals of length $t$, implies
$\|S(t)f-f\|_{\E}\to0$ as $t\downarrow0$.  The semigroup law yields strong
continuity at every time.

Positivity follows immediately from \eqref{eq:suspension}.  Finally, when
$t=n$ is an integer, $m(n,s)=n$ for almost every $s\in[0,1)$ and
\[
  (S(n)f)(s)=T^nf(s).
\]
The induced pointwise operator on $L^2([0,1);X)$ has norm $\|T^n\|$: the upper
bound is immediate, while the reverse inequality follows by testing constant
functions with values in an approximately norming unit vector for $T^n$.
\end{proof}

\begin{lemma}[Cesaro boundedness survives suspension]\label{lem:cesaro}
Assume
\begin{equation}\label{eq:disc-cesaro}
  C_T:=\sup_{N\geq1}\frac1N
  \left\|\sum_{k=0}^{N-1}T^k\right\|<\infty.
\end{equation}
Then the suspension \eqref{eq:suspension} is continuously Cesaro bounded:
\[
  \sup_{\tau>0}\frac1\tau
  \left\|\int_0^\tau S(t)\,dt\right\|<\infty.
\]
\end{lemma}

\begin{proof}
Put $B_N:=\sum_{k=0}^{N-1}T^k$.  From \eqref{eq:disc-cesaro},
\begin{equation}\label{eq:power-linear}
  \|B_N\|\leq C_TN,
  \qquad
  \|T^N\|=\|B_{N+1}-B_N\|\leq C_T(2N+1).
\end{equation}

Fix $N\geq1$ and $f\in\E$.  With Bochner integrals in $X$, set
\[
  g:=\int_0^1f(v)\,dv,
  \qquad h_-(s):=\int_s^1f(v)\,dv,
  \qquad h_+(s):=\int_0^sf(v)\,dv.
\]
Changing variables from $t$ to $s+t$ and splitting at the integer points gives,
for almost every $s\in[0,1)$,
\begin{equation}\label{eq:integer-integral}
  \int_0^N(S(t)f)(s)\,dt
  =h_-(s)+(B_N-I)g+T^Nh_+(s).
\end{equation}
Cauchy--Schwarz gives
\[
  \|g\|_X\leq\|f\|_{\E},
  \qquad \|h_-\|_{\E}\leq\|f\|_{\E},
  \qquad \|h_+\|_{\E}\leq\|f\|_{\E}.
\]
Consequently, \eqref{eq:power-linear} and \eqref{eq:integer-integral} imply
\[
  \left\|\int_0^NS(t)\,dt\right\|\leq C(1+N)
\]
with $C$ depending only on $C_T$.

For arbitrary $\tau=N+u$ with $N\in\N$ and $0\leq u<1$, the remaining
integral over $[N,N+u]$ has norm at most
\[
  \sup_{N\leq t\leq N+1}\|S(t)\|
  \leq\max\{\|T^N\|,\|T^{N+1}\|\}=O(N)
\]
by \eqref{eq:power-linear}.  Division by $\tau\geq N$ gives a uniform bound.
The range $0<\tau<1$ is controlled by local boundedness of the $C_0$-semigroup.
\end{proof}

\section{No uniform exponent gap}

\begin{theorem}\label{thm:main}
For every $\delta\in(0,1)$ there are an $L^2$-space $\E$ and a positive Kreiss
bounded $C_0$-semigroup $(S(t))_{t\geq0}$ on $\E$ such that
\begin{equation}\label{eq:lower}
  \|S(n)\|\geq\frac13(n+1)^{1-\delta}
  \qquad(n\in\N_0).
\end{equation}
Hence there is no number $\varepsilon_2>0$ for which every positive Kreiss
bounded $C_0$-semigroup on every $L^2$-space satisfies
$\|S(t)\|=O(t^{1-\varepsilon_2})$.
\end{theorem}

\begin{proof}
Fix $\delta\in(0,1)$.  By \cite[Theorems 2.2 and 2.3]{BonillaMuller}, there
are a Hilbert lattice $H$ and a positive Cesaro-bounded weighted shift
$T\in\mathcal L(H)$ such that
\[
  \|T^n\|\geq\frac13(n+1)^{1-\delta}
  \qquad(n\in\N_0).
\]
Their $H$ is a Hilbert direct sum of finite-dimensional coordinate lattices,
so $H$ is lattice isometric to an $\ell^2$-space.  Therefore
$\E=L^2([0,1);H)$ is an $L^2$-space.

Let $(S(t))$ be the suspension of $T$.  Lemma \ref{lem:basic} makes it a
positive $C_0$-semigroup, and Lemma \ref{lem:cesaro} makes it Cesaro bounded.
For an individually eventually positive semigroup on a Banach lattice,
continuous Cesaro boundedness is equivalent to Kreiss boundedness by
\cite[Proposition 2.6]{ArnoldCoine}; hence $(S(t))$ is Kreiss bounded.
Equation \eqref{eq:integer} now gives \eqref{eq:lower}.

If a universal $\varepsilon_2>0$ existed, choose
$0<\delta<\varepsilon_2$.  The universal upper bound at integer times would
contradict \eqref{eq:lower}, because
$n^{1-\delta}/n^{1-\varepsilon_2}=n^{\varepsilon_2-\delta}\to\infty$.
\end{proof}

\begin{remark}[Exact scope]
Theorem \ref{thm:main} does not disprove the literal source question.  For the
semigroup associated with a fixed $\delta$, it only forces every possible
polynomial gap $\varepsilon$ to satisfy $\varepsilon\leq\delta$.  It remains
open here whether each individual semigroup has some positive gap depending on
that semigroup.
\end{remark}

\section{Verification, limitations, and novelty}

No computation is used in the proof.  The two identities most relevant for
verification are \eqref{eq:integer-integral}, which transfers discrete Cesaro
boundedness to continuous time, and \eqref{eq:integer}, which transfers the
power lower bound without loss.  The former can also be checked directly by
partitioning the interval $[s,N+s]$ at its integer points.

A bounded literature check through 22 July 2026 covered
\cite{ArnoldCoine,BonillaMuller,CCEL,ArnoldCuny,Arora} and searches for
``positive Cesaro bounded power growth'', ``positive Kreiss bounded semigroup
sublinear polynomial'', and suspension variants.  No explicit statement of
Theorem \ref{thm:main} was found.  Novelty confidence is moderate-low: the
discrete example is known and the suspension is standard, while the exact
sharpness consequence for Question 3.9(2) appears not to have been recorded in
the checked sources.

The recommended human-review focus is Lemma \ref{lem:cesaro} and the invocation
of \cite[Proposition 2.6]{ArnoldCoine}.  Once those are accepted, the conclusion
is immediate from the Bonilla--Muller lower bound.

\begin{thebibliography}{9}

\bibitem{ArnoldCoine}
L. Arnold and C. Coine,
\emph{Growth rate of eventually positive Kreiss bounded $C_0$-semigroups on
$L^p$ and $\mathcal C(K)$},
J. Evol. Equ. 23 (2023), article 7;
\href{https://arxiv.org/abs/2207.08443}{arXiv:2207.08443}.

\bibitem{BonillaMuller}
A. Bonilla and V. Muller,
\emph{Kreiss bounded and uniformly Kreiss bounded operators},
\href{https://arxiv.org/abs/1912.07931}{arXiv:1912.07931}.

\bibitem{CCEL}
G. Cohen, C. Cuny, T. Eisner, and M. Lin,
\emph{Resolvent conditions and growth of powers of operators},
J. Math. Anal. Appl. 487 (2020), 124035;
\href{https://arxiv.org/abs/1912.10507}{arXiv:1912.10507}.

\bibitem{ArnoldCuny}
L. Arnold and C. Cuny,
\emph{On the growth rate of powers of a strongly Kreiss bounded operator on
an $L^p$-space},
\href{https://arxiv.org/abs/2302.14135}{arXiv:2302.14135}.

\bibitem{Arora}
S. Arora,
\emph{Eventually positive semigroups: spectral and asymptotic analysis},
Semigroup Forum 110 (2025), 263--295;
\href{https://arxiv.org/abs/2405.16371}{arXiv:2405.16371}.

\end{thebibliography}

\end{document}
